\documentclass[11pt]{article}
\usepackage[margin=2.1cm]{geometry}
\usepackage{booktabs}
\usepackage{amsmath}
\usepackage{enumitem}
\usepackage[T1]{fontenc}
\usepackage{lmodern}

\setlength{\parindent}{0pt}
\setlength{\parskip}{5pt}
\newcommand{\code}[1]{\texttt{#1}}
\newcommand{\answerbox}[1]{\vspace{2pt}\fbox{\parbox[c][#1][c]{\dimexpr\linewidth-2\fboxsep-2\fboxrule}{\ }}\vspace{4pt}}

\begin{document}

\begin{center}
{\Large\bfseries EE\,362 --- Pre-Lab 6}\\[2pt]
{\large Register File}\\[6pt]
\begin{tabular}{@{}p{6.5cm}p{4cm}p{4cm}@{}}
Name: \hrulefill & Section: Mon / Tue & Date: \hrulefill
\end{tabular}
\end{center}

\vspace{2pt}\hrule\vspace{4pt}

{\small\textbf{Due at the start of the lab session.} Five points, about twenty
minutes, on paper.}

\vspace{4pt}\hrule

\section*{Question 1 --- Paper first}

\textbf{(a)} Complete the write-behaviour table.

\begin{center}
\begin{tabular}{ccl@{\hskip 8pt}||@{\hskip 8pt}l}
\toprule
\# & \code{Strb} & \code{Dsel} & effect at the rising edge \\
\midrule
0 & 0 & --- & \rule{6cm}{0.4pt} \\[5pt]
1 & 1 & \code{"00000"} & \rule{6cm}{0.4pt} \\[5pt]
2 & 1 & $k \neq 0$ & \rule{6cm}{0.4pt} \\
\bottomrule
\end{tabular}
\end{center}

\textbf{(b)} Write the write-enable equation for register $k$, in terms of
\code{Strb}, the address minterm $m_k(\code{Dsel})$, and whatever else is needed:

\vspace{6pt}
$\mathit{We}_k = $ \rule{10cm}{0.4pt}
\vspace{8pt}

Circle the term that makes register 0 unwritable, and in one sentence say what
``hardwired to zero'' actually means in hardware.

\answerbox{1.6cm}

\textbf{(c)} A register file is three components you have already built. Name
them, and say which lab each came from.

\answerbox{2.0cm}

\section*{Question 2 --- Interface and cost}

\textbf{(a)} Write the VHDL \code{entity} for \code{RegFile}. Interface only.
Get the widths right: how many bits does an address need to select one of 32
registers?

\answerbox{4.4cm}

\textbf{(b)} Adding a \emph{third read port} costs very little. Adding a
\emph{second write port} is a much bigger problem. Explain the difference in two
or three sentences.

\answerbox{2.2cm}

\section*{Question 3 --- Test plan}

\begin{center}
\begin{tabular}{@{}c@{\hskip 8pt}p{5.0cm}@{\hskip 8pt}p{7.6cm}@{}}
\toprule
\# & what you do & what it proves \\
\midrule
1 & \rule{4.8cm}{0.4pt} & \rule{7.4cm}{0.4pt} \\[8pt]
2 & \rule{4.8cm}{0.4pt} & \rule{7.4cm}{0.4pt} \\[8pt]
3 & \rule{4.8cm}{0.4pt} & \rule{7.4cm}{0.4pt} \\[8pt]
4 & \rule{4.8cm}{0.4pt} & \rule{7.4cm}{0.4pt} \\[8pt]
5 & \rule{4.8cm}{0.4pt} & \rule{7.4cm}{0.4pt} \\[8pt]
6 & \rule{4.8cm}{0.4pt} & \rule{7.4cm}{0.4pt} \\
\bottomrule
\end{tabular}
\end{center}

\textbf{Then answer both:}

\textbf{(a)} Which of your vectors would catch a design that forgot to protect
register 0? If you do not have one, add it now.

\answerbox{1.4cm}

\textbf{(b)} Suppose a read port addresses the same register that is being
written on this clock edge. Two different designs could give two different
answers, and both would be correct designs. Name the two behaviours, and say how
your testbench would tell them apart.

\answerbox{2.4cm}

\vfill
{\small\emph{Question 3 is the one that matters.} It is the same question as
Section 5 of your lab report.}

\end{document}
